\section{Discussion}
\label{sec:discussion}

\subsection{What everyday rules afford}

The central claim of this work is that rules drawn from everyday life change what emergence can mean in an artwork. Three things support it.

First, everyday rules give interpretations something to be \emph{about}. A thought such as ``we walk by so many people'' can be tested against a system in which beings really do walk by each other, on their way between places they really do frequent. The same thought, read against Reas's elements or against a cellular automaton, would be a metaphor laid over an abstraction. Here it is a reading of a world that already contains walking, passing and staying.

Second, everyday rules make the system's emergent structure recognisable. The places in Section~\ref{sec:interpretations} are never drawn, yet three of the six interpretations reveal them, as the stars of a mesh, the bright fields of ripples and the centres of the separation's starbursts. This is not because the interpretations look for places, but because routines make beings return, and returning is what everyday life is mostly made of~\cite{gonzalez2008, song2010}. The interpretations converge on the structure of the system without being told it.

Third, everyday rules let a single system carry readings that are personal and still grounded. Each of the six thoughts comes from my own experience, and none of them was written to suit the system. That they can all be drawn from it suggests that a system made from life, even one as simple as this, holds more of life than its rules state, which is the promise of emergence in the first place.

\subsection{The interpretation as authorship}

In instruction art, the instruction belongs to the artist and its execution to someone else~\cite{lewitt1967, ono1964, obrist2013}. In this work both belong to me, yet the thought and the expression are still distinct, and the distance between them is where the work's judgements are made. Many interpretations went through several expressions before one matched its thought (the source tree keeps every earlier version), and in most cases it was the thought that held still.\note{check this against \code{wip/}: did the thoughts stay the same across versions while the expressions changed? if not, say what changed.} The three-part format is, in this sense, less a description of a finished work than a discipline for making one, and it is a format other artists can take up with other systems, or indeed with this one.

\subsection{Reading division into a neutral world}

\work{separation} deserves a separate comment, because it makes a claim that the system does not. Its thought speaks of beings that ``hold opposing views'', but beings in this system hold no views at all. Their colours are assigned at random, and they choose places without regard to anyone's colour. This is the opposite of Schelling's model of segregation, in which mild individual preferences about one's neighbours produce stark division~\cite{schelling1971}. In \work{separation}, there are no preferences, and the division is drawn entirely by the reader.

I think this is the interpretation's point rather than its weakness. The image looks like a map of a divided world, and it is made from beings who keep ordinary routines and have no opinions about each other. It asks how much of the division we see in maps of real places is in the people, and how much is in the reading.

\subsection{Limitations}

\paragraph{Not a model of a city.} The system is inspired by everyday life in New York, but it is not calibrated against any data, and it does not reproduce any measured property of real mobility. Places lie on a regular spiral, not a street network; a day is a year; there are no homes, no work, no weekends, no relationships, and no inheritance. Its rules are recognisable as life, not representative of it, and the claims made here are claims about an artwork.

\paragraph{Emergence, claimed qualitatively.} The emergent properties reported in Sections~\ref{sec:system} and~\ref{sec:revision} (a steady flow, neighbourhoods without homes, clustering with them) are measured, but the claims about what the interpretations reveal are made by looking. I have not measured, for instance, how closely the stars of the mesh coincide with places.

\paragraph{One artist.} All six interpretations are by the system's author. The claim that the format separates the maker of a system from its reader would be stronger if others read the same system; this is an invitation as much as a limitation (Section~\ref{sec:conclusion}).

\paragraph{Parameters as a confound.} Interpretations run in worlds of different sizes, densities and speeds. The rules are the same, but the worlds are not, and some differences between images come from their parameters rather than their readings.
